\chapter{Theories and Frameworks for Design}

When designing interactive systems, it is important to have a solid understanding of the theories and frameworks that underpin the design process. These theories provide a foundation for understanding user behavior, cognitive processes, and the principles of effective interaction design. In this chapter, we will explore several key theories and frameworks that are commonly used in the field of interactive systems design.

\section{Activity Theory}

\begin{figure}
    \centering
    \definecolor{activityBlue}{RGB}{68, 114, 196}
    \begin{tikzpicture}[
        % Configuración de los nodos circulares
        dot/.style={circle, fill=activityBlue, minimum size=5mm, inner sep=0pt},
        % Configuración de las etiquetas de texto
        labelStyle/.style={align=center},
        % Configuración de las flechas de doble punta
        arrow/.style={Stealth-Stealth, thick, activityBlue}
    ]
        % --- COORDENADAS DE LOS NODOS ---
        \node[dot] (rules) at (0,0) {};
        \node[dot] (community) at (4,0) {};
        \node[dot] (dol) at (8,0) {};
        \node[dot] (subject) at (2,2.8) {};
        \node[dot] (object) at (6,2.8) {};
        \node[dot] (instruments) at (4,5.6) {};

        % --- ETIQUETAS DE TEXTO ---
        \node[labelStyle, below=3mm] at (rules) {Rules};
        \node[labelStyle, below=3mm] at (community) {Community};
        \node[labelStyle, below=3mm] at (dol) {Division of Labour};
        \node[labelStyle, left=4mm] at (subject) {Subject};
        \node[labelStyle, right=4mm] at (object) {Object $\rightarrow$ Outcome};
        \node[labelStyle, above=3mm] at (instruments) {Instruments};

        % --- CONEXIONES (LÍNEAS Y FLECHAS) ---
        % Perímetro exterior del triángulo grande
        \draw[arrow] (subject) -- (rules);
        \draw[arrow] (instruments) -- (object);
        \draw[arrow] (rules) -- (community);
        \draw[arrow] (community) -- (dol);
        \draw[arrow] (dol) -- (object);

        % Conexiones internas desde el nodo superior (Instruments)
        \draw[arrow] (instruments) -- (subject);
        \draw[arrow] (instruments) -- (object);
        \draw[arrow] (instruments) -- (community);

        % Conexiones horizontales y cruzadas del medio
        \draw[arrow] (subject) -- (object);
        \draw[arrow] (subject) -- (community);
        \draw[arrow] (object) -- (community);
        
        % Conexiones cruzadas diagonales desde los extremos inferiores hacia el medio
        \draw[arrow] (rules) -- (object);
        \draw[arrow] (subject) -- (dol);
    \end{tikzpicture}
    \caption{Activity Theory Model}
    \label{fig:activity-theory}
\end{figure}

The Activity Theory model provides a comprehensive framework for understanding human activities in context. It provides three different hierarchical levels of acting:
\begin{enumerate}
    \item \ul{Activity}: This is the highest level of analysis and refers to the overall \emph{motive} that drives a person to engage in a particular action. Activities are typically long-term.
    \item \ul{Action}: This level focuses on the \emph{goal-directed} processes that individuals undertake to achieve specific goals. Actions are more immediate and are often part of a larger activity.
    \item \ul{Operation}: This is the most granular level of analysis and refers to the \emph{automatic} or habitual processes that individuals perform as part of an action. Operations are often unconscious.
\end{enumerate}

The Activity Theory model is represented in Figure \ref{fig:activity-theory}. It consists of several key components:
\begin{itemize}
    \item \ul{Subject}: The individual or group engaged in the activity. It interacts with the object through the use of instruments according to the rules and within the context of the community.
    \item \ul{Object}: The target and focus of the subject. It isn't neccessarily a physical object (e.g., questions) and its state usually changes as a result of the activity, so an outcome is produced.
    \item \ul{Instruments}: The tools or mediating artifacts that the subject uses to interact with and change the object. Instruments can be physical tools, symbols, or concepts.
    
    Depending on the context, an instrument can be an object in itself (as in the case of a software application, when the aim is creating it instead of using it).
    \ul{Rules}: The norms, conventions, and regulations that govern the activity.
    \item \ul{Community}: The group of individuals who share the same object and are involved in the activity. 
    \item \ul{Division of Labour}: The distribution of tasks and responsibilities among members of the community.
\end{itemize}

\section{Attribution Theory}

Attribution Theory is a psychological framework that seeks to explain how individuals interpret and assign causes to events and behaviors. It is particularly relevant in the context of interactive systems design, as it helps designers understand how users perceive and respond to system feedback, errors, and successes. It classifies attributions depending on different dimensions:
\begin{description}
    \item[Internal vs. External] This dimension distinguishes between causes that are attributed to the individual's own characteristics (internal) and those that are attributed to external factors, such as the environment or other people.
    \item[Stable vs. Variable] This dimension differentiates between causes that are perceived as stable and unchanging over time (stable) and those that are seen as temporary or fluctuating (variable).
    \item[Controllable vs. Uncontrollable] This dimension assesses whether the individual percieves the cause of an event as something they can influence or control (controllable) or as something beyond their control (uncontrollable).
\end{description}

In addition, the theory also considers the role of covariation in attributions. Covariation refers to the relationship between an effect and its potential causes over time. Three aspects are considered when analyzing covariation:
\begin{itemize}
    \item \ul{Consistency}: This aspect examines whether the attribution is similar to prior attributions made in similar situations. High consistency suggests that the cause is likely to be stable and reliable.
    \item \ul{Consensus}: This aspect looks at whether the attribution is similar to attributions made by others in similar situations. High consensus indicates that the cause is widely recognized and accepted.
    \item \ul{Distinctness}: This aspect assesses whether the attribution is directly associated with the specific situation. High distinctness suggests that the cause is unique to the particular context.
\end{itemize}

This theory is useful for designers to understand how users interpret system behavior and how they attribute success or failure to their own actions or to the system itself. If errors occure, describing them will help avoiding internal attributions, which can lead to frustration and negative user experiences.


\section{Motivation Theory}

Motivation Theory explores the factors that drive individuals to engage in certain behaviors and activities. In the context of interactive systems design, understanding user motivation is crucial for creating engaging and effective experiences. Two different approaches to motivation are commonly discussed: intrinsic and extrinsic motivation.
\begin{itemize}
    \item \ul{Intrinsic Motivation}: This type of motivation arises from within the individual and is driven by personal interest, enjoyment, or a sense of accomplishment. Users who are intrinsically motivated are more likely to engage with a system for the sheer pleasure of using it, rather than for external rewards or incentives.
    
    The \emph{Self-Determination Theory} (SDT) is a widely recognized framework that explains intrinsic motivation. According to SDT, individuals have three basic psychological needs that, when satisfied, enhance intrinsic motivation:
    \begin{enumerate}
        \item \ul{Competence}: The need to feel effective and capable in one's actions. When users perceive that they are mastering a task or skill, their intrinsic motivation is enhanced.
        \item \ul{Autonomy}: The need to feel a sense of control and choice over one's actions. When users have the freedom to make decisions and act according to their own preferences, their intrinsic motivation is strengthened.
        \item \ul{Relatedness}: The need to feel connected and supported by others. When users experience a sense of belonging and social connection, their intrinsic motivation is enhanced.
    \end{enumerate}
    \item \ul{Extrinsic Motivation}: This type of motivation is driven by external factors, such as rewards, recognition, or social pressure. Users who are extrinsically motivated may engage with a system to achieve specific outcomes or to gain approval from others, rather than for personal satisfaction.
\end{itemize}


\section{Place and Space Theory}

When designing interactive systems, it is important to consider the concepts of place and space, as they influence how users perceive and interact with their environment. An important distinction is made between the two concepts:
\begin{itemize}
    \item \ul{Space}: This refers to the physical environment that might enable or constrain user interactions. It encompasses the layout, dimensions, and spatial relationships of the environment in which users operate.
    \item \ul{Place}: This refers to the social and cultural context in which interactions occur. The space acquire meaning and significance through the activities, experiences, and social interactions that take place within it, resulting in a sense of place.
\end{itemize}

That way, the space is the opportunity for interaction, while the place is the context in which the interaction occurs. Designers must consider both aspects in order to achieve their created space to be perceived as a meaningful place by the users.